\section{Correctness}
\label{sec:correctness}

We separate what is tested, what is model-checked, and what is unresolved. Conflating the three is
how a systems paper overstates itself.

\subsection{Tested invariants}

\begin{enumerate}
  \item \textbf{No reuse without a passing guard pass.} Enforced by the typestate gate; the
  negative direction is exercised by deoptimisation tests across file content, metadata, symlink
  target, environment and input restoration.
  \item \textbf{No reuse of a corrupt object.} A corrupted CAS object produces a miss and a correct
  observation, measured directly.
  \item \textbf{No reuse for an ineligible class.} Clock, entropy, network, unmodelled syscall and
  writable shared mapping cases are fixtures in the integration suite and never hit.
  \item \textbf{No reuse from an empty cache}, and none after a cache relocation.
  \item \textbf{Enforcement is all-or-nothing.} A policy that cannot be installed produces an
  error, never an unenforced child labelled \proven{}; read and write rights are separate; an
  undeclared neighbour of a declared path is unreachable.
  \item \textbf{Restoration is verified before exposure.} Staged objects are hash-checked before
  any destination is mutated.
\end{enumerate}

Across every measurement cited in this paper the number of incorrect observations inside the
declared contract is zero.

\subsection{Model-checked properties}

The reuse state machine and the crash-recovery path are specified in TLA$^+\ \cite{lamport}$ and
checked with a bounded model checker over the hit, deopt, corrupt-object and restore-intent
transitions. The checked properties are the ordering ones: a hit is never emitted before the guard
pass, and a restore intent is never cleared before its outputs are durable. The model is a
specification of the design, not a verification of the Rust implementation, and we do not claim
otherwise.

\subsection{The residual counterexample}
\label{sec:toctou}

Guards are validated once, immediately before the workspace is mutated, and restoration happens
afterwards. A writer that changes an input \emph{inside that window} is not observed. We inject a
mutation at each point in the window and record the decision:

\begin{center}
\begin{tabular}{lll}
\toprule
Injection point & Decision & Observation \\
\midrule
before validation & deoptimised & correct for the new input \\
before loading cached streams & deoptimised & correct for the new input \\
mutate then revert & deoptimised & the mtime guard still detects it \\
\textbf{after validation, before apply} & \textbf{\texttt{Reused}} & \textbf{stale bytes} \\
concurrent writer during hashing (8/8) & deoptimised & correct for the new input \\
\bottomrule
\end{tabular}
\end{center}

This is a real counterexample, not a theoretical one: the run reports a hit and serves the previous
input's output. No economical userspace mechanism closes it. Descriptor pinning does not observe an
in-place rewrite through the same inode; a read-only bind mount does not freeze the inode against a
writer that already holds a writable description; snapshotting inputs redefines the observable
state and roughly doubles the hit cost; a cooperative lock is not a guarantee against an
uncooperative writer. Landlock confines the traced process --- it says nothing about unrelated
processes --- so enforcement does not narrow this window either.

We therefore state an \emph{operating assumption}: no external process mutates the inputs during
the reuse window. It is an assumption, not a guarantee, and the prerelease says so in its
installation instructions and in its safety document. TraceJIT must not be described as race-free,
atomic, or equivalent to fresh execution under concurrent mutation.

\subsection{Timestamp restoration}

A restored regular file carries its recorded metadata, including mtime. That is observable: a
program that prints its output's mtime sees the recorded value, and a build system that compares
mtimes can find a restored object older than an archive built from it, then skip re-archiving and
keep a stale member. We reproduced this on Lua 5.4.6 and zlib 1.3.1. Fresh timestamps were
measured and rejected: because standard output is replayed, a fresh mtime would make a replayed
observation describe a workspace that no longer exists.

\subsection{Enforcement lifecycle}

On a Landlock ABI~8 host, a read-only command completes the whole lifecycle: discovery
(\guarded{}), enforced execution (\proven{}), enforced reuse (\proven{}), mutation
(deoptimise, new bytes), re-execution (\proven{}). An output producer does not: the policy requires
its declared output to exist, and the command then recreates that output with
\texttt{O\_CREAT|O\_EXCL}, whose \texttt{EEXIST} is an unmodelled failed write. The run fails
closed. We tested whether a parent-directory replacement grant would unblock it; it does not, and
the grant was reverted rather than shipped. This is a limitation of the enforced tier, disclosed,
not a safety violation.
